% =====================================================================
%  2bat-ud1-5.tex  —  SESSIÓ 5: Combinar i escalar pressupostos
%  (sense preàmbul: s'inclou des de main.tex)
%  Criteris: C2 (sumar/restar i escalar amb interpretació). Estudi de cas.
% =====================================================================
\horatitol{Sessió 5 — Combinar i escalar pressupostos}%
{Un cas real: sumar dues partides d'un pressupost i aplicar-hi una retallada lineal amb el producte per un escalar.}

\subsection{Idea clau}

El pressupost d'una entitat sovint ve repartit en \emph{diverses partides}: les
despeses de personal en una taula i les de materials en una altra, totes dues
organitzades pels mateixos serveis i conceptes. Per saber quant costa cada cosa
\textbf{en total}, sumem les dues matrices (com a la sessió anterior).

Però els pressupostos també canvien de manera \emph{global}: un finançador
anuncia una retallada lineal del $10\%$, o un conveni aplica una pujada del $5\%$
a tot. Aquests canvis «a tot alhora» es fan amb el \textbf{producte d'una matriu
per un escalar}: es multiplica \emph{cada} casella pel mateix nombre.

\begin{idea}
\textbf{Dues operacions, dues lectures socials}
\begin{itemize}[leftmargin=1.4em, itemsep=2pt]
  \item \textbf{Suma de matrices:} integra dues partides (personal $+$ materials).
        Cal que tinguin la mateixa dimensió i es fa casella a casella.
  \item \textbf{Producte per un escalar:} aplica un percentatge global. Per a una
        retallada del $10\%$ es multiplica per $0{,}9$ (queda el $90\%$); per a un
        increment del $5\%$, per $1{,}05$.
  \item \textbf{Combinar-les:} primer es pot sumar i després escalar, per exemple
        $0{,}9\cdot(A+B)$: s'obté la despesa total ja retallada.
\end{itemize}
\end{idea}

\subsection{Exemples}

\begin{exemple}[sumar personal i materials]
Una entitat reparteix el pressupost mensual de dos serveis (files) en dues
columnes (Servei 1 i Servei 2). Les despeses de personal i de materials són
\[
\text{Personal: } A=\begin{pmatrix} 1800 & 1200 \\ 1500 & 900 \end{pmatrix},\qquad
\text{Materials: } B=\begin{pmatrix} 400 & 300 \\ 350 & 250 \end{pmatrix}.
\]
La despesa total per servei i concepte és la suma:
\[
A+B=\begin{pmatrix} 1800+400 & 1200+300 \\ 1500+350 & 900+250 \end{pmatrix}
    =\begin{pmatrix} 2200 & 1500 \\ 1850 & 1150 \end{pmatrix}.
\]
\end{exemple}

\begin{exemple}[aplicar una retallada del 10\%]
Sobre la despesa total $T=\begin{pmatrix} 2200 & 1500 \\ 1850 & 1150 \end{pmatrix}$
el finançador imposa una retallada lineal del $10\%$. El nou pressupost és el
$90\%$, és a dir, multipliquem per l'escalar $0{,}9$:
\[
0{,}9\cdot T=\begin{pmatrix} 0{,}9\per 2200 & 0{,}9\per 1500 \\ 0{,}9\per 1850 & 0{,}9\per 1150 \end{pmatrix}
            =\begin{pmatrix} 1980 & 1350 \\ 1665 & 1035 \end{pmatrix}.
\]
La primera casella passa de $2200\eur$ a $1980\eur$: la retallada hi treu $220\eur$.
La retallada lineal afecta \emph{tots} els conceptes amb el mateix percentatge.
\end{exemple}

\subsection{Exercicis resolts}

\begin{exresolt}[la despesa total i la retallada]
Una associació té el pressupost trimestral de tres serveis (files) repartit en
personal i materials (columnes):
\[
P=\begin{pmatrix} 2000 & 600 \\ 1600 & 500 \\ 1200 & 400 \end{pmatrix},\qquad
M=\begin{pmatrix} 500 & 200 \\ 400 & 100 \\ 300 & 200 \end{pmatrix}.
\]
\begin{enumerate}[label=\alph*), leftmargin=1.6em, topsep=2pt]
  \item Troba la despesa total $P+M$.
  \item S'aplica una retallada lineal del $10\%$. Quin escalar cal? Troba el nou
        pressupost.
  \item Quant s'estalvia, en total, el primer servei (fila 1) amb la retallada?
\end{enumerate}
\solucio
a) Sumem casella a casella (totes dues són $3\times 2$):
\[
P+M=\begin{pmatrix} 2500 & 800 \\ 2000 & 600 \\ 1500 & 600 \end{pmatrix}.
\]
b) Una retallada del $10\%$ deixa el $90\%$: multipliquem per $0{,}9$:
\[
0{,}9\cdot(P+M)=\begin{pmatrix} 2250 & 720 \\ 1800 & 540 \\ 1350 & 540 \end{pmatrix}.
\]
c) El primer servei tenia $2500+800=3300\eur$ i ara té $2250+720=2970\eur$:
s'estalvia $3300-2970=330\eur$ (que és el $10\%$ de $3300$).
\resposta{$P+M=\begin{pmatrix} 2500 & 800 \\ 2000 & 600 \\ 1500 & 600 \end{pmatrix}$; \;amb $0{,}9$: $\begin{pmatrix} 2250 & 720 \\ 1800 & 540 \\ 1350 & 540 \end{pmatrix}$; \;el primer servei estalvia $330\eur$.}
\end{exresolt}

\begin{exresolt}[un increment en lloc d'una retallada]
El pressupost anual de dos casals (files) en dos conceptes (columnes) és
\[
C=\begin{pmatrix} 4000 & 1000 \\ 3000 & 800 \end{pmatrix}.
\]
Un nou conveni aplica un increment lineal del $5\%$ a tot el pressupost. Per quin
escalar cal multiplicar? Troba el nou pressupost i digues quant puja, en total, el
segon casal.
\solucio
Un increment del $5\%$ vol dir que es paga el $105\%$: multipliquem per $1{,}05$:
\[
1{,}05\cdot C=\begin{pmatrix} 1{,}05\per 4000 & 1{,}05\per 1000 \\ 1{,}05\per 3000 & 1{,}05\per 800 \end{pmatrix}
            =\begin{pmatrix} 4200 & 1050 \\ 3150 & 840 \end{pmatrix}.
\]
El segon casal tenia $3000+800=3800\eur$ i ara té $3150+840=3990\eur$: puja
$3990-3800=190\eur$ (el $5\%$ de $3800$).
\resposta{$1{,}05\cdot C=\begin{pmatrix} 4200 & 1050 \\ 3150 & 840 \end{pmatrix}$; \;el segon casal puja $190\eur$.}
\end{exresolt}

\subsection{Exercicis per fer}

\begin{exercicis}
\begin{exlist}
  \item El pressupost mensual de dos serveis (files) en personal i materials
        (columnes) és
        \[ A=\begin{pmatrix} 1500 & 400 \\ 1200 & 300 \end{pmatrix},\qquad
           B=\begin{pmatrix} 300 & 100 \\ 200 & 150 \end{pmatrix}. \]
        Troba la despesa total $A+B$.
  \item A partir de la despesa total $T=\begin{pmatrix} 1800 & 500 \\ 1400 & 450 \end{pmatrix}$,
        aplica una retallada lineal del $10\%$ (multiplica per $0{,}9$) i troba el
        nou pressupost.
  \item Un pressupost és $D=\begin{pmatrix} 2000 & 800 \\ 1000 & 600 \end{pmatrix}$.
        \begin{enumerate}[label=\alph*), leftmargin=1.6em, topsep=2pt]
          \item Aplica-hi una pujada del $5\%$ (escalar $1{,}05$).
          \item Per quin escalar multiplicaries per a una retallada del $20\%$?
        \end{enumerate}
  \item Una entitat suma personal i materials i després aplica una retallada:
        \[ P=\begin{pmatrix} 2500 & 700 \\ 1500 & 500 \end{pmatrix},\qquad
           M=\begin{pmatrix} 500 & 300 \\ 500 & 200 \end{pmatrix}. \]
        Calcula $0{,}9\cdot(P+M)$ (primer suma, després retalla el $10\%$).
  \item Després d'aplicar una retallada del $10\%$, un servei ha quedat amb un
        pressupost de $1800\eur$. Quin era el pressupost \emph{abans} de la
        retallada? (Pista: $1800$ és el $90\%$ del pressupost original.)
  \item Quin escalar cal per a cadascun d'aquests canvis lineals?
        \begin{enumerate}[label=\alph*), leftmargin=1.6em, topsep=2pt]
          \item una retallada del $15\%$ \quad b) un increment del $10\%$ \quad
                c) reduir el pressupost a la meitat
        \end{enumerate}
  \item \textbf{(Raonament)} Una companya diu que «retallar primer un $10\%$ i
        després sumar els materials» dona el mateix que «sumar primer els materials
        i després retallar el $10\%$». Té raó? Comprova-ho amb
        $A=\begin{pmatrix} 1000 & 500 \end{pmatrix}$ i
        $B=\begin{pmatrix} 200 & 100 \end{pmatrix}$, i explica què passa amb la
        retallada quan només es retalla el personal.
\end{exlist}
\end{exercicis}

% ---------------------------------------------------------------------
\fullsolucions{Sessió 5}
\begin{solucions}
\begin{sollist}
  \item $A+B=\begin{pmatrix} 1800 & 500 \\ 1400 & 450 \end{pmatrix}$.
  \item $0{,}9\cdot T=\begin{pmatrix} 1620 & 450 \\ 1260 & 405 \end{pmatrix}$.
  \item (a) $1{,}05\cdot D=\begin{pmatrix} 2100 & 840 \\ 1050 & 630 \end{pmatrix}$
        \quad (b) per l'escalar $0{,}8$ (es manté el $80\%$).
  \item $P+M=\begin{pmatrix} 3000 & 1000 \\ 2000 & 700 \end{pmatrix}$;
        \;$0{,}9\cdot(P+M)=\begin{pmatrix} 2700 & 900 \\ 1800 & 630 \end{pmatrix}$.
  \item El pressupost original era $1800:0{,}9=2000\eur$.
  \item (a) $0{,}85$ \quad (b) $1{,}1$ \quad (c) $0{,}5$.
  \item Sí, dona el mateix: $0{,}9\cdot(A+B)=0{,}9\cdot\begin{pmatrix} 1200 & 600 \end{pmatrix}
        =\begin{pmatrix} 1080 & 540 \end{pmatrix}$, igual que
        $0{,}9\cdot A + 0{,}9\cdot B=\begin{pmatrix} 900 & 450 \end{pmatrix}+\begin{pmatrix} 180 & 90 \end{pmatrix}$.
        En canvi, si la retallada s'aplica \emph{només} al personal i després s'hi
        sumen els materials sencers, el resultat és diferent: l'ordre i el que es
        retalla importen.
\end{sollist}
\end{solucions}
